% Einheit 2 – Satz vom Nullprodukt
\clearpage
\setcounter{aufgabe}{18}
\hypertarget{e2}{}\einheitenkopf{Einheit 2 von 4 · Satz vom Nullprodukt}
\zweigzeile{Hier lernst du, Gleichungen wie $(x - a)\cdot(x - b) = 0$ mit dem Satz vom Nullprodukt zu lösen, auch durch Ausklammern von $x$ · neu oder schon bekannt – je nach Buch (an der Oberschule erst in Klasse 10) · P10 · baut auf: lineare Gleichungen, Ausklammern, Einsetzen prüfen}

\begin{aufgabe}{Ich erkenne, ob bei einem Produkt rechts eine Null steht. Steht auf einer Seite ein Produkt und auf der anderen null, gilt der Satz vom Nullprodukt, sonst nicht. Kreuze an. Löse nichts.}
\begin{teile}
\teil $(x - 1)\cdot(x - 5) = 0$ \hfill \kreuz{gilt}\kreuz{gilt nicht}
\teil $x\cdot(x + 4) = 21$ \hfill \kreuz{gilt}\kreuz{gilt nicht}
\teil $(x + 2)\cdot(x - 7) = 0$ \hfill \kreuz{gilt}\kreuz{gilt nicht}
\teil $(x - 3)\cdot x = -2$ \hfill \kreuz{gilt}\kreuz{gilt nicht}
\teil $0 = x\cdot(x - 1)$ \hfill \kreuz{gilt}\kreuz{gilt nicht}
\end{teile}
\end{aufgabe}

\begin{aufgabe}{Ich kann eine Gleichung in Produktform mit dem Satz vom Nullprodukt lösen. Setze jeden Faktor null.}
\beispiel{(x - 2)\cdot(x - 5) &= 0 \\ x - 2 &= 0 \quad\text{oder}\quad x - 5 = 0 \\ x_1 &= 2, \quad x_2 = 5}
\begin{gleichungsraster}[2]
\gl{(x - 1)\cdot(x - 4) = 0} & \gl{(x - 3)\cdot(x - 6) = 0} \\ \noalign{\vspace{6pt}}
\gl{(x - 8)\cdot(x - 2) = 0} & \gl{(x - 5)\cdot(x - 7) = 0} \\ \noalign{\vspace{6pt}}
\gl{(x + 4)\cdot(x - 1) = 0} & \gl{x\cdot(x + 6) = 0} \\ \noalign{\vspace{6pt}}
\gl{(x + 3)\cdot(x + 3) = 0} & \\ \noalign{\vspace{6pt}}
\end{gleichungsraster}
\medskip
\begin{teile}
\teil Prüfe durch Einsetzen, ob $x = -2$ eine Lösung von $(x - 3)\cdot(x + 1) = 0$ ist. \hfill \kreuz{(wA)}\kreuz{(fA)}
\end{teile}
\end{aufgabe}

\begin{aufgabe}{Ich kann eine Gleichung lösen, indem ich $x$ ausklammere. Klammere $x$ aus und setze dann jeden Faktor null.}
\beispiel{x^2 + 3x &= 0 \\ x\cdot(x + 3) &= 0 \\ x_1 &= 0, \quad x_2 = -3}
\begin{gleichungsraster}[2]
\gl{x^2 + 4x = 0} & \gl{x^2 + 9x = 0} \\ \noalign{\vspace{6pt}}
\gl{x^2 - 8x = 0} & \gl{2x^2 + 10x = 0} \\ \noalign{\vspace{6pt}}
\end{gleichungsraster}
\medskip
Rechts steht keine Null: Multipliziere aus und bringe alles auf eine Seite. Lösen kannst du die Gleichung in Einheit 3.
\begin{gleichungsraster}[2]
\gl{x\cdot(x + 1) = 30} & \\ \noalign{\vspace{6pt}}
\end{gleichungsraster}
\medskip
Klammere zuerst die Zahl vor $x^2$ aus. Schreibe dann den Rest als Produkt zweier Klammern und löse.
\begin{gleichungsraster}[3]
\gl{2x^2 - 2x - 24 = 0} & \\ \noalign{\vspace{6pt}}
\end{gleichungsraster}
\medskip
\begin{teile}
\teil Welcher Wert ist eine Lösung der Gleichung $x\cdot(x + 7) = -12$? Kreuze an. (P10 2025)\\ \kreuz{$3$}\kreuz{$4$}\kreuz{$-3$}\kreuz{$-12$}
\end{teile}
\end{aufgabe}

\begin{aufgabe}{Ich finde den Fehler, wenn jemand durch $x$ teilt. Jonas löst $x^2 = 6x$ so:}
\rechnung{x^2 &= 6x &&\mid :x \\ x &= 6}
\begin{teile}
\teil Finde den Fehler, benenne ihn und korrigiere die Rechnung.
\end{teile}
\medskip
Löse selbst.
\begin{gleichungsraster}[3]
\gl{4x^2 = 28x} & \\ \noalign{\vspace{6pt}}
\end{gleichungsraster}
\end{aufgabe}

\begin{aufgabe}{Ich kann begründen, wann der Satz vom Nullprodukt gilt.}
\begin{teile}
\teil Begründe ohne Rechnung, warum $(x - 4)\cdot(x + 1) = 0$ die Lösungen $4$ und $-1$ hat.
\teil Ben löst $x\cdot(x - 1) = 6$ so: „$x = 6$ oder $x - 1 = 6$.“ Hat Ben recht? Begründe.
\end{teile}
\end{aufgabe}

\begin{aufgabe}{Ich kann mit dem Satz vom Nullprodukt die Spannweite eines Brückenbogens bestimmen. Ein Brückenbogen hat die Form der Parabel zu $y = -0{,}05\cdot x\cdot(x - 48)$. Dabei sind $x$ und $y$ in Metern angegeben, der Boden liegt bei $y = 0$.}
\begin{teile}
\teil An welchen Stellen steht der Bogen auf dem Boden? Wie groß ist die Spannweite des Bogens?
\teil Der höchste Punkt liegt genau in der Mitte. Wie hoch ist der Bogen dort?
\end{teile}
\end{aufgabe}
